\NormalFont\ShortTitle{1 Petru}

{\MT Prima scrisoare a lui Petru


\par }\ChapOne{1}{\PP \VerseOne{1}Petru, apostol al lui Isus Hristos, către cei aleși, care locuiesc ca străini în Dispersia, în Pont, în Galatia, în Capadocia, în Asia și în Bitinia,

\VS{2}după știința dinainte a lui Dumnezeu Tatăl, în sfințirea Duhului, ca să vă supuneți lui Isus Hristos și să fiți stropiți cu sângele Lui: Harul să vă fie vouă și pacea să se înmulțească.


\par }{\PP \VS{3}Binecuvântat să fie Dumnezeul și Tatăl Domnului nostru Isus Hristos, care, după marea Sa îndurare, ne-a făcut să ne renaștem la o nădejde vie, prin învierea lui Isus Hristos din morți,

\VS{4}la o moștenire incoruptibilă și nestricăcioasă, care nu se pierde, rezervată în ceruri pentru voi,

\VS{5}care, prin puterea lui Dumnezeu, sunteți păziți prin credință pentru o mântuire care se va arăta în vremea de apoi.

\VS{6}În aceasta vă bucurați foarte mult, deși acum, pentru puțină vreme, dacă este nevoie, ați fost îndurerați în diferite încercări,

\VS{7}pentru ca dovada credinței voastre, care este mai prețioasă decât aurul care piere, chiar dacă este încercat prin foc, să se dovedească a avea ca rezultat lauda, gloria și onoarea la arătarea lui Isus Hristos —

\VS{8}pe care, fără să-L cunoașteți, Îl iubiți. În el, deși acum nu-l vedeți, crezând, vă bucurați foarte mult, cu o bucurie de nedescris și plină de glorie,

\VS{9}primind rezultatul credinței voastre, mântuirea sufletelor voastre.


\par }{\PP \VS{10}Proorocii au căutat și au cercetat cu sârguință această mântuire. Ei au profețit despre harul care va veni la voi,

\VS{11}căutând la cine sau la ce fel de timp arăta Duhul lui Hristos care era în ei când a prezis suferințele lui Hristos și gloria care le va urma.

\VS{12}Lor li s-a arătat că nu ei își slujeau lor înșiși, ci vouă, în aceste lucruri care v-au fost anunțate acum prin cei care v-au vestit Buna Vestire, prin Duhul Sfânt trimis din ceruri; lucruri pe care îngerii doresc să le cerceteze.


\par }{\PP \VS{13}De aceea, pregătiți-vă mințile pentru acțiune. Fiți cumpătați și puneți-vă toată nădejdea în harul care vă va fi adus la descoperirea lui Isus Hristos—

\VS{14}ca niște copii ai ascultării, fără să vă conformați după poftele voastre de odinioară, ca în ignoranța voastră,

\VS{15}ci, după cum cel care v-a chemat este sfânt, fiți și voi sfinți în toată purtarea voastră,

\VS{16}pentru că este scris: “Să fiți sfinți, căci eu sunt sfânt”.


\par }{\PP \VS{17}Dacă Îl chemați ca Tată pe El, care, fără deosebire de persoane, judecă după fapta fiecăruia, petreceți aici, ca străini, timpul vieții voastre, cu frică și respect,

\VS{18}știind că ați fost răscumpărați, nu cu lucruri coruptibile, ca argintul sau aurul, de la modul de viață nefolositor, moștenit de la părinții voștri,

\VS{19}ci cu sânge prețios, ca al unui miel fără cusur și fără pată, sângele lui Hristos,

\VS{20}care a fost cunoscut mai dinainte de întemeierea lumii, dar care a fost descoperit în acest ultim veac pentru voi,

\VS{21}care, prin el, sunteți credincioși în Dumnezeu, care l-a înviat din morți și i-a dat slavă, pentru ca credința și speranța voastră să fie în Dumnezeu.


\par }{\PP \VS{22}După ce v-ați curățit sufletele în ascultarea adevărului prin Duhul Sfânt, într-o sinceră afecțiune frățească, iubiți-vă cu ardoare unii pe alții din inimă,

\VS{23}fiind născuți din nou, nu din sămânță coruptibilă, ci din incoruptibilă, prin cuvântul lui Dumnezeu, care trăiește și rămâne în veci.

\VS{24}Căci,

\par }{\Q “Toată carnea este ca iarba,

\par }{\QB și toată slava omului ca floarea din iarbă.

\par }{\Q Iarba se usucă, iar floarea ei cade;


\par }{\QB \VS{25}dar cuvântul Domnului este veșnic.”

\par }{\PP Acesta este cuvântul de Buna Vestire care v-a fost propovăduit.



\par }\Chap{2}{\PP \VerseOne{1}Deci, lepădând orice răutate, orice înșelăciune, fățărnicie, invidie și orice vorbire de rău,

\VS{2}ca niște prunci nou-născuți, tânjiți după laptele curat duhovnicesc, ca să creșteți cu el,

\VS{3}dacă într-adevăr ați gustat că Domnul este milostiv.

\VS{4}Veniți la el, piatră vie, respinsă într-adevăr de oameni, dar aleasă de Dumnezeu, prețioasă.

\VS{5}Și voi, ca pietre vii, sunteți zidiți ca o casă spirituală, pentru a fi o preoție sfântă, care să aducă jertfe spirituale, plăcute lui Dumnezeu prin Isus Cristos.

\VS{6}Pentru că așa este cuprins în Scriptură,

\par }{\Q “Iată, pun în Sion o piatră de temelie, aleasă și prețioasă.

\par }{\QB Cel care crede în El nu va fi dezamăgit.”


\par }{\PP \VS{7}Pentru voi, cei ce credeți, este deci cinstea, dar pentru cei neascultători,

\par }{\Q “Piatra pe care au respins-o zidarii

\par }{\QB a devenit piatra de temelie,”


\par }{\PP \VS{8}și,

\par }{\Q “o piatră de poticnire și o stâncă de jignire”.

\par }{\PP Căci ei se poticnesc de cuvânt, fiind neascultători, la care au și fost rânduiți.

\VS{9}Dar voi sunteți un neam ales, o preoție regală, o națiune sfântă, un popor care este proprietatea lui Dumnezeu, ca să vestiți măreția celui care v-a chemat din întuneric la lumina Lui minunată.

\VS{10}În trecut, voi nu erați un popor, dar acum sunteți poporul lui Dumnezeu, care nu avuseseră milă, dar acum ați obținut milă.


\par }{\PP \VS{11}Preaiubiților, vă rog, ca străini și pelerini, să vă feriți de poftele trupești, care luptă împotriva sufletului,

\VS{12}și să vă purtați bine printre neamuri, pentru ca, atunci când vă vor vorbi de rău, ca de niște răufăcători, să vadă faptele voastre bune și să slăvească pe Dumnezeu în ziua vizitei.


\par }{\PP \VS{13}Supuneți-vă deci, pentru Domnul, oricărei legi omenești, fie împăratului, ca stăpânitor,

\VS{14}fie guvernatorilor, ca trimiși de El pentru a se răzbuna pe cei răi și pentru a lăuda pe cei ce fac bine.

\VS{15}Căci aceasta este voia lui Dumnezeu: ca, prin fapte bune, să puneți capăt ignoranței oamenilor nebuni.

\VS{16}Trăiți ca niște oameni liberi, dar fără să vă folosiți libertatea ca pe o mantie a răutății, ci ca niște robi ai lui Dumnezeu.


\par }{\PP \VS{17}Cinstiți pe toți oamenii. Iubiți frăția. Temeți-vă de Dumnezeu. Onorați-l pe rege.


\par }{\PP \VS{18}Robilor, supuneți-vă stăpânilor voștri cu toată stima, nu numai celor buni și blânzi, ci și celor răi.

\VS{19}Căci este lăudabil dacă cineva suportă durerea, suferind pe nedrept, din cauza conștiinței față de Dumnezeu.

\VS{20}Căci ce glorie este aceasta dacă, atunci când păcătuiești, înduri cu răbdare bătăile? Dar dacă, atunci când faci bine, înduri cu răbdare suferința, acest lucru este lăudabil înaintea lui Dumnezeu.

\VS{21}Căci la aceasta ați fost chemați, pentru că și Cristos a suferit pentru noi, lăsându-vă o pildă, ca să urmați pașii lui,

\VS{22}care nu a păcătuit, “și nu s-a găsit înșelăciune în gura lui”.

\VS{23}Când a fost blestemat, nu a răspuns la blestem. Când a suferit, nu a amenințat, ci s-a încredințat celui care judecă cu dreptate.

\VS{24}El însuși a purtat păcatele noastre în trupul Său pe lemn, pentru ca noi, după ce am murit față de păcate, să trăim pentru dreptate. Voi ați fost vindecați prin rănile Lui.

\VS{25}Căci voi vă rătăceați ca niște oi, dar acum v-ați întors la Păstorul și Supraveghetorul sufletelor voastre.



\par }\Chap{3}{\PP \VerseOne{1}De asemenea, soțiilor, supuneți-vă soților voștri, pentru ca, chiar dacă unii nu ascultă de Cuvânt, să fie câștigați prin purtarea soțiilor lor, fără să le spună nimic,

\VS{2}văzând purtarea voastră curată și cu frică.

\VS{3}Frumusețea voastră să nu vină din podoaba exterioară, din împletirea părului, din purtarea de podoabe de aur sau din îmbrăcarea cu haine frumoase,

\VS{4}ci din persoana ascunsă a inimii, în podoaba incoruptibilă a unui spirit blând și liniștit, care este foarte prețioasă în ochii lui Dumnezeu.

\VS{5}Căci așa se împodobeau în trecut și femeile sfinte care sperau în Dumnezeu, supunându-se soților lor.

\VS{6}Astfel, Sara a ascultat de Avraam, numindu-l stăpân, ai cărui copii sunteți acum, dacă faceți bine și nu vă înspăimântați de nicio spaimă.


\par }{\PP \VS{7}Voi, soților, trăiți la fel cu soțiile voastre, în cunoștință de cauză, dând cinste femeii ca unui vas mai slab, ca și cum ați fi împreună-moștenitori ai harului vieții, pentru ca rugăciunile voastre să nu fie împiedicate.


\par }{\PP \VS{8}În sfârșit, fiți cu toții la fel, milostivi, iubitori ca niște frați, blânzi la suflet, politicoși,

\VS{9}fără să dați rău pentru rău și fără să dați insultă pentru insultă, ci binecuvântând, știind că ați fost chemați la aceasta, ca să moșteniți binecuvântarea.

\VS{10}Căci,

\par }{\Q “Cel care vrea să iubească viața

\par }{\QB și să vezi zile bune,

\par }{\Q să-și țină limba departe de rău

\par }{\QB și buzele lui să nu mai spună minciuni.


\par }{\Q \VS{11}Să se abată de la rău și să facă binele.

\par }{\QB Să caute pacea și să o urmărească.


\par }{\Q \VS{12}Căci ochii Domnului sunt asupra celor drepți,

\par }{\QB și urechile sale deschise la rugăciunile lor;

\par }{\QB dar fața Domnului este împotriva celor ce fac răul.”


\par }{\PP \VS{13}Și cine vă va face rău, dacă veți fi imitatori ai binelui?

\VS{14}Dar chiar dacă veți suferi din pricina dreptății, sunteți binecuvântați. “Nu vă temeți de ceea ce se tem ei și nu vă tulburați.”

\VS{15}Ci sfințiți pe Domnul Dumnezeu în inimile voastre. Fiți întotdeauna gata să dați un răspuns oricui vă întreabă un motiv cu privire la speranța care este în voi, cu umilință și teamă,

\VS{16}având o conștiință bună. Astfel, în timp ce se vorbește despre voi ca despre niște răufăcători, pot fi dezamăgiți cei care vă blestemă modul vostru bun de viață în Cristos.

\VS{17}Căci este mai bine, dacă este voia lui Dumnezeu, să suferiți pentru că faceți ceea ce este drept decât pentru că faceți răul.

\VS{18}Fiindcă și Hristos a suferit o dată pentru păcate, cel drept pentru cei nedrepți, ca să vă aducă la Dumnezeu, fiind omorât în trup, dar înviat în Duhul,

\VS{19}în care a mers și el și a predicat duhurilor din închisoare,

\VS{20}care mai înainte au fost neascultători, când Dumnezeu a așteptat cu răbdare, în zilele lui Noe, în timp ce se construia corabia. În ea, puțini, adică opt suflete, au fost salvați prin apă.

\VS{21}Acesta este un simbol al botezului, care acum vă mântuiește — nu îndepărtarea murdăriei cărnii, ci răspunsul unei conștiințe bune față de Dumnezeu — prin învierea lui Isus Cristos,

\VS{22}care se află la dreapta lui Dumnezeu, după ce a urcat în ceruri, îngerii, autoritățile și puterile fiindu-i supuse.



\par }\Chap{4}{\PP \VerseOne{1}Deci, de vreme ce Hristos a suferit pentru noi în trup, înarmați-vă și voi cu același gând; căci cel ce a suferit în trup a încetat de la păcat,

\VS{2}ca să nu mai trăiți în trup pentru poftele oamenilor, ci pentru voia lui Dumnezeu.

\VS{3}Căci am petrecut destul din timpul trecut făcând poftele neamurilor și umblând în desfrâu, în pofte, în beții, în orgii, în desfrâuri și în idolatrii abominabile.

\VS{4}Lor li se pare ciudat că voi nu alergați cu ei în același exces de dezmăț, vorbind de rău despre voi.

\VS{5}Ei vor da socoteală celui care este gata să judece pe cei vii și pe cei morți.

\VS{6}Căci în acest scop a fost propovăduită Buna Vestire chiar și morților, pentru ca ei să fie judecați într-adevăr ca oameni în carne și oase, dar să trăiască ca pentru Dumnezeu în duh.


\par }{\PP \VS{7}Dar sfârșitul tuturor lucrurilor este aproape. De aceea, fiți sănătoși la minte, stăpâni pe voi înșivă și sobri în rugăciune.

\VS{8}Și, mai presus de toate, fiți sârguincioși în dragostea voastră între voi, căci dragostea acoperă o mulțime de păcate.

\VS{9}Fiți ospitalieri unii cu alții, fără să vă certați.

\VS{10}După cum fiecare a primit un dar, folosiți-l în slujirea reciprocă, ca buni administratori ai harului lui Dumnezeu în diferitele sale forme.

\VS{11}Dacă cineva vorbește, să fie ca și cum ar fi chiar cuvintele lui Dumnezeu. Dacă cineva slujește, să fie ca din puterea pe care o dă Dumnezeu, pentru ca în toate lucrurile să fie glorificat Dumnezeu prin Isus Hristos, căruia îi aparțin gloria și stăpânirea în vecii vecilor. Amin.


\par }{\PP \VS{12}Preaiubiților, nu vă mirați de încercarea de foc care a venit peste voi ca să vă încerce, ca și cum vi s-ar fi întâmplat ceva ciudat.

\VS{13}Ci, pentru că sunteți părtași la suferințele lui Hristos, bucurați-vă, pentru ca, la arătarea gloriei Lui, să vă bucurați și voi cu bucurie nespus de mare.

\VS{14}Dacă sunteți insultați pentru numele lui Hristos, sunteți binecuvântați, pentru că Duhul slavei și al lui Dumnezeu se odihnește peste voi. Din partea lor, el este hulit, dar din partea voastră este glorificat.

\VS{15}Dar niciunul dintre voi să nu sufere ca un ucigaș, sau ca un hoț, sau ca un făcător de rele, sau ca un băgător de seamă în treburile altora.

\VS{16}Dar dacă unul dintre voi suferă pentru că este creștin, să nu se rușineze, ci să slăvească pe Dumnezeu în această privință.

\VS{17}Căci a venit vremea ca judecata să înceapă cu casa lui Dumnezeu. Dacă începe mai întâi cu noi, ce se va întâmpla cu cei care nu ascultă de Vestea Bună a lui Dumnezeu?

\VS{18}“Dacă este greu pentru cel neprihănit să fie salvat, ce se va întâmpla cu cel nelegiuit și cu păcătosul?”

\VS{19}De aceea, și cei care suferă după voia lui Dumnezeu, făcând binele, să îi încredințeze sufletele lor, ca unui Creator credincios.



\par }\Chap{5}{\PP \VerseOne{1}De aceea îndemn pe bătrânii din mijlocul vostru, ca un bătrân și ca martor al suferințelor lui Hristos, care va avea parte și el de slava care se va arăta:

\VS{2}Păstoriți turma lui Dumnezeu care este în mijlocul vostru, exercitând supravegherea, nu de bunăvoie, ci de bunăvoie; nu pentru un câștig necinstit, ci de bunăvoie;

\VS{3}nu ca un stăpân peste cei care v-au fost încredințați, ci făcându-vă voi înșivă exemple pentru turmă.

\VS{4}Când se va arăta Păstorul cel mare, voi veți primi cununa gloriei care nu se ofilește.


\par }{\PP \VS{5}Și voi, cei mai tineri, fiți supuși celui mai mare. Da, îmbrăcați-vă cu toții cu smerenie și supuneți-vă unii altora; căci “Dumnezeu se împotrivește celor mândri, dar dă har celor smeriți”.

\VS{6}Umiliți-vă, așadar, sub mâna puternică a lui Dumnezeu, pentru ca el să vă înalțe la timpul potrivit,

\VS{7}aruncând asupra lui toate grijile voastre, pentru că el are grijă de voi.


\par }{\PP \VS{8}Fii treaz și stăpân pe tine însuți. Fiți vigilenți. Adversarul vostru, diavolul, umblă ca un leu care răcnește, căutând pe cine să devoreze.

\VS{9}Rezistați-i cu tărie în credința voastră, știind că frații voștri care sunt în lume trec prin aceleași suferințe.

\VS{10}Dar fie ca Dumnezeul oricărui har, care v-a chemat la gloria Sa eternă prin Hristos Isus, după ce ați suferit puțin, să vă desăvârșească, să vă întărească, să vă întărească și să vă așeze.

\VS{11}A lui să fie gloria și puterea în vecii vecilor. Amin.


\par }{\PP \VS{12}Prin Silvanus, fratele nostru credincios, așa cum îl consider, v-am scris pe scurt, îndemnându-vă și mărturisind că acesta este adevăratul har al lui Dumnezeu în care vă aflați.

\VS{13}Cea care este în Babilon, aleasă împreună cu voi, vă salută. La fel și Marcu, fiul meu.

\VS{14}Salutați-vă unii pe alții cu un sărut de dragoste.

\par }{\PP Pacea fie cu toți cei care sunteți în Hristos Isus. Amin.


\par }